\begin{theorem}[Chebyshev’s sum inequality]
Let  
\[
a_{1}\le a_{2}\le\cdots\le a_{n},\qquad 
b_{1}\le b_{2}\le\cdots\le b_{n}
\]
be real numbers. Then  
\[
n\sum_{i=1}^{n}a_i b_i\;\ge\;
\Bigl(\sum_{i=1}^{n}a_i\Bigr)
\Bigl(\sum_{i=1}^{n}b_i\Bigr),
\tag{1}
\]
with equality iff at least one of the two sequences is constant
($a_{1}=a_{2}= \dots =a_{n}$ or $b_{1}=b_{2}= \dots =b_{n}$).
\end{theorem}

\begin{proof}
\textbf{1. Degenerate cases.}  
If $n=1$ both sides of \eqref{1} coincide, so the inequality holds with equality.  
If all $a_i$ are equal, say $a_i\equiv c$, then the left–hand side of \eqref{1} is $nc\sum b_i$ and the right–hand side is $c\sum b_i\cdot n$, which are identical; the same argument applies when all $b_i$ are equal. Hence the statement is true in these degenerate situations.  
In the sequel we assume $n\ge2$ and that neither sequence is constant.

\medskip\noindent\textbf{2. Algebraic reformulation.}  
Set  
\[
\bar a:=\frac1n\sum_{i=1}^{n}a_i,\qquad 
\bar b:=\frac1n\sum_{i=1}^{n}b_i .
\]
Then
\begin{align*}
n\sum_{i=1}^{n}a_i b_i-\Bigl(\sum_{i=1}^{n}a_i\Bigr)\Bigl(\sum_{i=1}^{n}b_i\Bigr)
&=n\sum_{i=1}^{n}a_i b_i-n^{2}\bar a\,\bar b   \\
&=n\Bigl[\sum_{i=1}^{n}a_i b_i-n\bar a\,\bar b\Bigr].
\end{align*}
Moreover
\[
\sum_{i=1}^{n}(a_i-\bar a)(b_i-\bar b)
   =\sum_{i=1}^{n}a_i b_i-n\bar a\,\bar b .
\]
Consequently
\begin{equation}
n\sum_{i=1}^{n}a_i b_i-\Bigl(\sum_{i=1}^{n}a_i\Bigr)\Bigl(\sum_{i=1}^{n}b_i\Bigr)
   =n\sum_{i=1}^{n}(a_i-\bar a)(b_i-\bar b).
\tag{2}
\end{equation}
Since $n>0$, inequality \eqref{1} is equivalent to
\begin{equation}
\sum_{i=1}^{n}(a_i-\bar a)(b_i-\bar b)\ge0 .
\tag{3}
\end{equation}
Define the centred sequences
\[
\alpha_i:=a_i-\bar a,\qquad \beta_i:=b_i-\bar b\qquad(i=1,\dots ,n).
\]
Subtracting a constant does not change the order, so
\[
\alpha_1\le\alpha_2\le\cdots\le\alpha_n,\qquad 
\beta_1\le\beta_2\le\cdots\le\beta_n .
\]
Thus $(\alpha_i)$ and $(\beta_i)$ are **similarly ordered**.

\medskip\noindent\textbf{3. Application of the rearrangement inequality.}  
The rearrangement inequality states that for any permutation $\sigma$ of $\{1,\dots ,n\}$,
\begin{equation}
\sum_{i=1}^{n}\alpha_i\beta_i
   \;\ge\;\frac1{n!}\sum_{\sigma}\sum_{i=1}^{n}\alpha_i\beta_{\sigma(i)} .
\tag{4}
\end{equation}
Because the centred sequences have zero sum,
\[
\sum_{i=1}^{n}\alpha_i=0=\sum_{i=1}^{n}\beta_i,
\]
the inner double sum equals $0$, and therefore the right–hand side of \eqref{4} is $0$. Hence
\[
\sum_{i=1}^{n}\alpha_i\beta_i\ge0 .
\]
Using the definitions of $\alpha_i,\beta_i$ this is precisely \eqref{3}, and by \eqref{2} we obtain \eqref{1}.

\medskip\noindent\textbf{4. Equality condition.}  
Equality in the rearrangement inequality occurs iff the two sequences are equal up to a permutation or one of them is constant.  
Since $\sum\alpha_i=0=\sum\beta_i$, the “equal up to a permutation’’ case forces every $\alpha_i$ and $\beta_i$ to be zero; otherwise the sum $\sum\alpha_i\beta_i$ could not be zero. Hence equality in \eqref{3} (and thus in \eqref{1}) holds precisely when

\begin{itemize}
\item all $\alpha_i=0$, i.e. $a_1=a_2=\dots =a_n$, or
\item all $\beta_i=0$, i.e. $b_1=b_2=\dots =b_n$.
\end{itemize}
These are exactly the cases stated in the theorem (the case $n=1$ being included).

\qed
\end{proof}